% Generated draft: requires mathematical and source review.
\documentclass[11pt]{article}
\usepackage[margin=25mm]{geometry}
\usepackage{fontspec}
\setmainfont{Latin Modern Roman}
\usepackage{amsmath,amssymb,mathtools,unicode-math}
\setmathfont{Latin Modern Math}
\usepackage{enumitem,needspace,xurl,hyperref}
\setlength{\parindent}{0pt}
\setlength{\parskip}{5pt}
\setlength{\emergencystretch}{3em}
\newcommand{\N}{\mathbb{N}}
\newcommand{\Z}{\mathbb{Z}}
\newcommand{\Q}{\mathbb{Q}}
\newcommand{\R}{\mathbb{R}}
\newcommand{\C}{\mathbb{C}}
\newcommand{\F}{\mathbb{F}}
\newcommand{\A}{\mathbb{A}}
\newcommand{\PP}{\mathbb{P}}
\newcommand{\E}{\mathbb{E}}
\newcommand{\eps}{\varepsilon}
\newcommand{\dd}{\,\mathrm{d}}
\begin{document}
\section*{UnsolvedMath ID: 20000347}
\textbf{Source problem code:} catalogue entry
\subsection{Definitions and mathematical statement}
Grushevsky asks (1) to construct divisors of small slope on $\mathcal A_g$, and (2) to find explicit complete subvarieties of $\mathcal A_g$, noting a codimension bound in [KS03]. Here $\mathcal A_g$ is the moduli space or stack of principally polarized abelian varieties of dimension $g$, with dimension $g(g+1)/2$. Completeness means properness inside the open moduli space, not merely in a compactification. The source does not specify a target dimension, numerical slope bound, field, stack/coarse convention or compactification.

To define slope, choose a compactification with Hodge class $\lambda$ and specified boundary class $\delta$. For an effective divisor closure of numerical class $a\lambda-b\delta$ with $b>0$, its slope is $s=a/b$ under this convention. The treatment of multiple boundary components and normalization of $\delta$ must be fixed; a slope is not intrinsic to an unspecified divisor class on the open space alone.

The attached unreviewed aid treats $g\ge4$. It extracts the non-two-torsion component $N'_0$ from Mumford's cycle relation
\[
N_0=\theta_{\mathrm{null}}+2N'_0,
\]
where the Andreotti--Mayer divisor $N_0$ parametrizes singular theta divisors and $\theta_{\mathrm{null}}$ concerns singular even two-torsion points. It claims an effective class $a_g\lambda-b_g\delta$ with $b_g>0$ satisfying
\[
6<s(N'_0)<6+\frac{12}{g+1}\qquad(g\ge4).
\]
The aid reports exact coefficient and slope-defect formulas but does not include them; they cannot be recovered from the inequalities alone.

For a complete $d$-dimensional subvariety $X\subset\mathcal A_g$, the aid reports the boundary-restriction observation
\[
\delta|_X=0,\qquad
(a\lambda-b\delta)|_X=a\lambda|_X,\qquad
\int_X(a\lambda-b\delta)^d=a^d\int_X\lambda^d.
\]
Thus restriction and internal intersection on such an interior cycle do not detect $b$ or the ratio $a/b$. This limitation does not construct complete subvarieties or establish their maximal dimension.

The aid's literature assessment says exact minimum effective slopes are known only through $g=5$ in its checked sources, and mentions a submitted November 2025 preprint claiming the exact maximum dimension of complete subvarieties, without identifying it or its formula. Those are supplied dated assessments, not an independent current literature check. Both original construction requests remain broader than the stated slope example and restriction observation.
\subsection{Short English statement}
Construct low-slope effective divisors and explicit complete subvarieties of principally polarized abelian moduli. The aid gives a slope bound for one divisor and explains why intersections wholly inside the open moduli space cannot measure its boundary coefficient.
\subsection{Sources}
\begin{itemize}
\item Grushevsky, Problem 6.1 in AIM \emph{Compact moduli spaces and birational geometry}, ``Moduli of abelian varieties'': \url{https://aimath.org/WWN/birational/birational.pdf}. Canonical location: \texttt{aim-algebraic-geometry-notes.json}, notes[346]. The supplied verification resolves the split section/problem number; no PDF was consulted.
\item Original reference [KS03], for a codimension bound on complete subvarieties. Its authors, title and precise bound are not supplied.
\item The attached aid names Mumford's divisor-cycle relation and mentions an unidentified submitted November 2025 preprint. Full citations, exact coefficient formulas and the claimed maximal-dimension formula are absent. The aid is machine generated, not peer reviewed, with low novelty confidence.
\item Local registry: \texttt{data/unsolvedmath/problems.json}, permanent catalogue entry, exact source code \texttt{catalogue entry}.
\end{itemize}
\end{document}
